% !TEX program = xelatex
% Bu dosya tex_uret.py tarafından sorular.py'den üretilmiştir; düzenleme için sorular.py'yi değiştirip yeniden üretin.
\documentclass[11pt,a4paper]{article}
\usepackage[margin=1.8cm,top=1.9cm,bottom=1.9cm]{geometry}
\usepackage{fontspec}
\usepackage{polyglossia}
\setdefaultlanguage{turkish}
\usepackage{amsmath,amssymb}
\usepackage{graphicx}
\usepackage{booktabs}
\usepackage{array}
\usepackage{enumitem}
\usepackage{multicol}
\usepackage[most]{tcolorbox}
\usepackage{fancyhdr}
\usepackage{titlesec}
\usepackage{xcolor}
\usepackage{needspace}
\usepackage[hidelinks]{hyperref}

\definecolor{anamavi}{RGB}{0,45,98}
\definecolor{cozumyesil}{RGB}{0,110,50}
\definecolor{uyarikirmizi}{RGB}{180,20,20}

\titleformat{\section}{\large\bfseries\color{anamavi}}{\thesection.}{0.5em}{}
\titleformat{\subsection}{\normalsize\bfseries\color{anamavi}}{\thesubsection}{0.5em}{}

\pagestyle{fancy}\fancyhf{}
\lhead{\small MMAT105 Genel Matematik}
\rhead{\small Hafta 15 -- Final Sınavı Deneme}
\cfoot{\small\thepage}
\renewcommand{\headrulewidth}{0.4pt}

\newtcolorbox{kural}[1][Kural]{colback=anamavi!5,colframe=anamavi,title=#1,fonttitle=\bfseries,breakable,left=4pt,right=4pt,top=2pt,bottom=2pt}
\newtcolorbox{dikkat}[1][Dikkat]{colback=uyarikirmizi!5,colframe=uyarikirmizi,title=#1,fonttitle=\bfseries,breakable,left=4pt,right=4pt,top=2pt,bottom=2pt}
\newenvironment{cozum}{\begin{tcolorbox}[colback=cozumyesil!6,colframe=cozumyesil,boxrule=0.5pt,left=4pt,right=4pt,top=2pt,bottom=2pt]\textbf{\color{cozumyesil}Çözüm.}\ }{\end{tcolorbox}}
\newcommand{\sk}[1]{\textbf{\color{anamavi}#1)}~}
\setlength{\parindent}{0pt}
\setlength{\parskip}{3pt}
\setlength{\lineskiplimit}{2pt}\setlength{\lineskip}{4pt}
\setlength{\columnsep}{0.9cm}
\setlength{\columnseprule}{0.4pt}
\newlength{\soruara}\setlength{\soruara}{2.3cm}

\begin{document}

\begin{minipage}[c]{\textwidth}
{\LARGE\bfseries\color{anamavi} MMAT105 Genel Matematik}\\[4pt]
{\large\bfseries Final Sınavı Deneme}\\[2pt]
{\normalsize Hafta 1--14 konuları}\\[6pt]
{\small 2026--2027 Güz}
\end{minipage}

\vspace{4pt}\noindent\textcolor{anamavi}{\rule{\textwidth}{1pt}}

\begin{minipage}[t]{0.55\textwidth}
\vspace{2pt}
\textbf{Ad Soyad:} \rule{0.68\linewidth}{0.4pt}\\[12pt]
\textbf{Öğrenci No:} \rule{0.6\linewidth}{0.4pt}
\end{minipage}\hfill
\begin{minipage}[t]{0.42\textwidth}
\begin{kural}[Yönerge]
\small\raggedright
\begin{itemize}[leftmargin=*,itemsep=0pt,topsep=0pt]
\item Soru sayısı: \(25\); \ önerilen süre: \textbf{\(75\) dakika}.
\item Her soru eşit puanlıdır (\(4\) puan).
\item Her sorunun yalnızca \textbf{bir} doğru cevabı vardır.
\item Yanlış cevaplar doğru cevapları götürmez.
\end{itemize}
\end{kural}
\end{minipage}

\vspace{6pt}
\begin{multicols}{2}
\noindent\begin{minipage}{\linewidth}\raggedright
\textbf{\color{anamavi}1.}\ \((-2)^3-(-3)^2+4\cdot 5-12\div 3\) işleminin sonucu kaçtır?
\par\smallskip\noindent{\setlength{\tabcolsep}{1pt}\begin{tabular*}{\linewidth}{@{}l@{\extracolsep{\fill}}llll@{}}\sk{A}\(-17\) & \sk{B}\(-9\) & \sk{C}\(-1\) & \sk{D}\(1\) & \sk{E}\(17\)\end{tabular*}}
\end{minipage}
\par\vspace{\soruara}
\noindent\begin{minipage}{\linewidth}\raggedright
\textbf{\color{anamavi}2.}\ \(84\) sayısının farklı asal bölenlerinin toplamı kaçtır?
\par\smallskip\noindent{\setlength{\tabcolsep}{1pt}\begin{tabular*}{\linewidth}{@{}l@{\extracolsep{\fill}}llll@{}}\sk{A}\(7\) & \sk{B}\(12\) & \sk{C}\(14\) & \sk{D}\(16\) & \sk{E}\(21\)\end{tabular*}}
\end{minipage}
\par\vspace{\soruara}
\noindent\begin{minipage}{\linewidth}\raggedright
\textbf{\color{anamavi}3.}\ Bir duraktan üç otobüs hattının otobüsleri sırasıyla \(15\), \(20\) ve \(25\) dakikada bir kalkmaktadır. Üç hattın otobüsleri saat 08.00'de birlikte kalktığına göre, üçü yeniden ilk kez birlikte saat kaçta kalkar?
\par\smallskip\noindent{\setlength{\tabcolsep}{1pt}\begin{tabular*}{\linewidth}{@{}l@{\extracolsep{\fill}}llll@{}}\sk{A}09.00 & \sk{B}09.15 & \sk{C}09.40 & \sk{D}12.00 & \sk{E}13.00\end{tabular*}}
\end{minipage}
\par\vspace{\soruara}
\noindent\begin{minipage}{\linewidth}\raggedright
\textbf{\color{anamavi}4.}\ \(2-\cfrac{1}{1-\cfrac{1}{3}}\) işleminin sonucu kaçtır?
\par\smallskip\noindent{\setlength{\tabcolsep}{1pt}\begin{tabular*}{\linewidth}{@{}l@{\extracolsep{\fill}}llll@{}}\sk{A}\(\dfrac{1}{2}\) & \sk{B}\(\dfrac{2}{3}\) & \sk{C}\(\dfrac{4}{3}\) & \sk{D}\(\dfrac{3}{2}\) & \sk{E}\(\dfrac{5}{2}\)\end{tabular*}}
\end{minipage}
\par\vspace{\soruara}
\noindent\begin{minipage}{\linewidth}\raggedright
\textbf{\color{anamavi}5.}\ Ali'nin parası Ayşe'nin parasının \(3\) katıdır. Ali, Ayşe'ye \(40\) TL verirse paraları eşit oluyor. Ali'nin başlangıçtaki parası kaç TL'dir?
\par\smallskip\noindent{\setlength{\tabcolsep}{1pt}\begin{tabular*}{\linewidth}{@{}l@{\extracolsep{\fill}}llll@{}}\sk{A}\(40\) & \sk{B}\(80\) & \sk{C}\(100\) & \sk{D}\(120\) & \sk{E}\(160\)\end{tabular*}}
\end{minipage}
\par\vspace{\soruara}
\noindent\begin{minipage}{\linewidth}\raggedright
\textbf{\color{anamavi}6.}\ \(30\) soruluk bir sınavda her doğru cevap \(4\) puan kazandırmakta, her yanlış cevap \(1\) puan kaybettirmektedir. Tüm soruları cevaplayan bir öğrenci \(70\) puan aldığına göre kaç soruya doğru cevap vermiştir?
\par\smallskip\noindent{\setlength{\tabcolsep}{1pt}\begin{tabular*}{\linewidth}{@{}l@{\extracolsep{\fill}}llll@{}}\sk{A}\(14\) & \sk{B}\(18\) & \sk{C}\(20\) & \sk{D}\(22\) & \sk{E}\(25\)\end{tabular*}}
\end{minipage}
\par\vspace{\soruara}
\noindent\begin{minipage}{\linewidth}\raggedright
\textbf{\color{anamavi}7.}\ \(2<|x-1|\le 5\) eşitsizliğini sağlayan kaç tane \(x\) tam sayısı vardır?
\par\smallskip\noindent{\setlength{\tabcolsep}{1pt}\begin{tabular*}{\linewidth}{@{}l@{\extracolsep{\fill}}llll@{}}\sk{A}\(4\) & \sk{B}\(6\) & \sk{C}\(7\) & \sk{D}\(8\) & \sk{E}\(10\)\end{tabular*}}
\end{minipage}
\par\vspace{\soruara}
\noindent\begin{minipage}{\linewidth}\raggedright
\textbf{\color{anamavi}8.}\ \(\dfrac{2^5\cdot 4^3}{8^4}\) işleminin sonucu kaçtır?
\par\smallskip\noindent{\setlength{\tabcolsep}{1pt}\begin{tabular*}{\linewidth}{@{}l@{\extracolsep{\fill}}llll@{}}\sk{A}\(\dfrac{1}{2}\) & \sk{B}\(2\) & \sk{C}\(4\) & \sk{D}\(8\) & \sk{E}\(16\)\end{tabular*}}
\end{minipage}
\par\vspace{\soruara}
\noindent\begin{minipage}{\linewidth}\raggedright
\textbf{\color{anamavi}9.}\ \(x\) bir gerçek sayı olmak üzere \(\dfrac{2^{x+4}-2^{x}}{2^{x+1}+2^{x}}\) ifadesinin değeri kaçtır?
\par\smallskip\noindent{\setlength{\tabcolsep}{1pt}\begin{tabular*}{\linewidth}{@{}l@{\extracolsep{\fill}}llll@{}}\sk{A}\(3\) & \sk{B}\(4\) & \sk{C}\(5\) & \sk{D}\(\dfrac{16}{3}\) & \sk{E}\(8\)\end{tabular*}}
\end{minipage}
\par\vspace{\soruara}
\noindent\begin{minipage}{\linewidth}\raggedright
\textbf{\color{anamavi}10.}\ \(a=2^{48},\quad b=3^{36},\quad c=5^{24}\) sayılarının doğru sıralanışı aşağıdakilerden hangisidir?
\par\smallskip\noindent\begin{tabular}{@{}p{0.47\linewidth}p{0.47\linewidth}@{}}\sk{A}\(a<b<c\) & \sk{B}\(b<c<a\)\\[7pt] \sk{C}\(c<a<b\) & \sk{D}\(a<c<b\)\\[7pt] \sk{E}\(c<b<a\) & \end{tabular}
\end{minipage}
\par\vspace{\soruara}
\noindent\begin{minipage}{\linewidth}\raggedright
\textbf{\color{anamavi}11.}\ \(\sqrt{98}+\sqrt{8}-\sqrt{32}\) işleminin sonucu kaçtır?
\par\smallskip\noindent\begin{tabular}{@{}p{0.47\linewidth}p{0.47\linewidth}@{}}\sk{A}\(\sqrt{74}\) & \sk{B}\(5\sqrt{2}\)\\[7pt] \sk{C}\(9\sqrt{2}\) & \sk{D}\(13\sqrt{2}\)\\[7pt] \sk{E}\(3\sqrt{2}\) & \end{tabular}
\end{minipage}
\par\vspace{\soruara}
\noindent\begin{minipage}{\linewidth}\raggedright
\textbf{\color{anamavi}12.}\ \(\dfrac{6}{\sqrt{7}-1}-\sqrt{7}\) işleminin sonucu kaçtır?
\par\smallskip\noindent{\setlength{\tabcolsep}{1pt}\begin{tabular*}{\linewidth}{@{}l@{\extracolsep{\fill}}llll@{}}\sk{A}\(1\) & \sk{B}\(\sqrt{7}-1\) & \sk{C}\(2\) & \sk{D}\(\sqrt{7}+1\) & \sk{E}\(6\)\end{tabular*}}
\end{minipage}
\par\vspace{\soruara}
\noindent\begin{minipage}{\linewidth}\raggedright
\textbf{\color{anamavi}13.}\ \(\sqrt{11+\sqrt{72}}\) ifadesinin eşiti aşağıdakilerden hangisidir?
\par\smallskip\noindent\begin{tabular}{@{}p{0.47\linewidth}p{0.47\linewidth}@{}}\sk{A}\(\sqrt{6}+\sqrt{5}\) & \sk{B}\(2\sqrt{2}+\sqrt{3}\)\\[7pt] \sk{C}\(2+\sqrt{7}\) & \sk{D}\(3+2\sqrt{2}\)\\[7pt] \sk{E}\(3+\sqrt{2}\) & \end{tabular}
\end{minipage}
\par\vspace{\soruara}
\noindent\begin{minipage}{\linewidth}\raggedright
\textbf{\color{anamavi}14.}\ \(x^2-y^2=48\) ve \(x+y=12\) olduğuna göre \(x\cdot y\) kaçtır?
\par\smallskip\noindent{\setlength{\tabcolsep}{1pt}\begin{tabular*}{\linewidth}{@{}l@{\extracolsep{\fill}}llll@{}}\sk{A}\(4\) & \sk{B}\(16\) & \sk{C}\(20\) & \sk{D}\(24\) & \sk{E}\(32\)\end{tabular*}}
\end{minipage}
\par\vspace{\soruara}
\noindent\begin{minipage}{\linewidth}\raggedright
\textbf{\color{anamavi}15.}\ \(a+b=10\) ve \(a\cdot b=21\) olduğuna göre \(a^2+b^2\) kaçtır?
\par\smallskip\noindent{\setlength{\tabcolsep}{1pt}\begin{tabular*}{\linewidth}{@{}l@{\extracolsep{\fill}}llll@{}}\sk{A}\(16\) & \sk{B}\(58\) & \sk{C}\(79\) & \sk{D}\(100\) & \sk{E}\(142\)\end{tabular*}}
\end{minipage}
\par\vspace{\soruara}
\noindent\begin{minipage}{\linewidth}\raggedright
\textbf{\color{anamavi}16.}\ \(x\notin\{-3,\,-2,\,0,\,3\}\) olmak üzere \(\dfrac{x^2-x-6}{x^2-9}\div\dfrac{x+2}{x}\) ifadesinin en sade biçimi aşağıdakilerden hangisidir?
\par\smallskip\noindent\begin{tabular}{@{}p{0.47\linewidth}p{0.47\linewidth}@{}}\sk{A}\(\dfrac{x+2}{x+3}\) & \sk{B}\(\dfrac{x}{x-3}\)\\[7pt] \sk{C}\(\dfrac{x-3}{x}\) & \sk{D}\(\dfrac{x}{x+3}\)\\[7pt] \sk{E}\(\dfrac{(x+2)^2}{x(x+3)}\) & \end{tabular}
\end{minipage}
\par\vspace{\soruara}
\noindent\begin{minipage}{\linewidth}\raggedright
\textbf{\color{anamavi}17.}\ \(\dfrac{a}{b}=\dfrac{5}{8}\) ve \(a+b=91\) olduğuna göre \(b\) kaçtır?
\par\smallskip\noindent{\setlength{\tabcolsep}{1pt}\begin{tabular*}{\linewidth}{@{}l@{\extracolsep{\fill}}llll@{}}\sk{A}\(21\) & \sk{B}\(35\) & \sk{C}\(49\) & \sk{D}\(56\) & \sk{E}\(63\)\end{tabular*}}
\end{minipage}
\par\vspace{\soruara}
\noindent\begin{minipage}{\linewidth}\raggedright
\textbf{\color{anamavi}18.}\ \(8\) bahçıvan günde \(6\) saat çalışarak bir parkın bakımını \(5\) günde bitirmektedir. Aynı hızda çalışan \(6\) bahçıvan günde \(4\) saat çalışarak aynı işi kaç günde bitirir?
\par\smallskip\noindent{\setlength{\tabcolsep}{1pt}\begin{tabular*}{\linewidth}{@{}l@{\extracolsep{\fill}}llll@{}}\sk{A}\(10\) & \sk{B}\(12\) & \sk{C}\(15\) & \sk{D}\(16\) & \sk{E}\(20\)\end{tabular*}}
\end{minipage}
\par\vspace{\soruara}
\noindent\begin{minipage}{\linewidth}\raggedright
\textbf{\color{anamavi}19.}\ Bir sınıftaki \(24\) öğrencinin not ortalaması \(65\)'tir. Sınıfa, not ortalaması \(80\) olan \(k\) öğrenci katıldığında sınıfın not ortalaması \(70\) oluyor. Buna göre \(k\) kaçtır?
\par\smallskip\noindent{\setlength{\tabcolsep}{1pt}\begin{tabular*}{\linewidth}{@{}l@{\extracolsep{\fill}}llll@{}}\sk{A}\(4\) & \sk{B}\(6\) & \sk{C}\(8\) & \sk{D}\(10\) & \sk{E}\(12\)\end{tabular*}}
\end{minipage}
\par\vspace{\soruara}
\noindent\begin{minipage}{\linewidth}\raggedright
\textbf{\color{anamavi}20.}\ Hangi sayının \(\dfrac{2}{5}\)'inin \(3\) fazlası \(15\)'tir?
\par\smallskip\noindent{\setlength{\tabcolsep}{1pt}\begin{tabular*}{\linewidth}{@{}l@{\extracolsep{\fill}}llll@{}}\sk{A}\(18\) & \sk{B}\(24\) & \sk{C}\(30\) & \sk{D}\(36\) & \sk{E}\(45\)\end{tabular*}}
\end{minipage}
\par\vspace{\soruara}
\noindent\begin{minipage}{\linewidth}\raggedright
\textbf{\color{anamavi}21.}\ Bir annenin yaşı ile kızının yaşının toplamı bugün \(50\)'dir. \(4\) yıl önce annenin yaşı, kızının yaşının \(5\) katıydı. Annenin bugünkü yaşı kaçtır?
\par\smallskip\noindent{\setlength{\tabcolsep}{1pt}\begin{tabular*}{\linewidth}{@{}l@{\extracolsep{\fill}}llll@{}}\sk{A}\(35\) & \sk{B}\(39\) & \sk{C}\(41\) & \sk{D}\(45\) & \sk{E}\(46\)\end{tabular*}}
\end{minipage}
\par\vspace{\soruara}
\noindent\begin{minipage}{\linewidth}\raggedright
\textbf{\color{anamavi}22.}\ Bir öğrenci harçlığının \(\dfrac{1}{3}\)'ü ile kitap, kalan paranın \(\dfrac{1}{4}\)'ü ile defter almıştır. Geriye \(300\) TL kaldığına göre öğrencinin harçlığı kaç TL'dir?
\par\smallskip\noindent{\setlength{\tabcolsep}{1pt}\begin{tabular*}{\linewidth}{@{}l@{\extracolsep{\fill}}llll@{}}\sk{A}\(300\) & \sk{B}\(360\) & \sk{C}\(400\) & \sk{D}\(450\) & \sk{E}\(600\)\end{tabular*}}
\end{minipage}
\par\vspace{\soruara}
\noindent\begin{minipage}{\linewidth}\raggedright
\textbf{\color{anamavi}23.}\ \(40\,000\) TL, yıllık \(\%24\) basit faiz oranıyla \(9\) ay süreyle bankaya yatırılıyor. Bu süre sonunda elde edilen faiz kaç TL'dir?
\par\smallskip\noindent{\setlength{\tabcolsep}{1pt}\begin{tabular*}{\linewidth}{@{}l@{\extracolsep{\fill}}llll@{}}\sk{A}\(800\) & \sk{B}\(4800\) & \sk{C}\(7200\) & \sk{D}\(9600\) & \sk{E}\(47\,200\)\end{tabular*}}
\end{minipage}
\par\vspace{\soruara}
\noindent\begin{minipage}{\linewidth}\raggedright
\textbf{\color{anamavi}24.}\ Bir satıcı, maliyeti \(1250\) TL olan bir ürünün etiket fiyatını maliyetin \(\%40\) fazlası olarak belirliyor. Ürünü etiket fiyatı üzerinden \(\%20\) indirimle satan satıcının kârı kaç TL'dir?
\par\smallskip\noindent{\setlength{\tabcolsep}{1pt}\begin{tabular*}{\linewidth}{@{}l@{\extracolsep{\fill}}llll@{}}\sk{A}\(150\) & \sk{B}\(250\) & \sk{C}\(280\) & \sk{D}\(350\) & \sk{E}\(500\)\end{tabular*}}
\end{minipage}
\par\vspace{\soruara}
\noindent\begin{minipage}{\linewidth}\raggedright
\textbf{\color{anamavi}25.}\ Tuz oranı \(\%20\) olan \(300\) gram tuzlu su ile tuz oranı \(\%35\) olan kaç gram tuzlu su karıştırılırsa tuz oranı \(\%30\) olan bir karışım elde edilir?
\par\smallskip\noindent{\setlength{\tabcolsep}{1pt}\begin{tabular*}{\linewidth}{@{}l@{\extracolsep{\fill}}llll@{}}\sk{A}\(150\) & \sk{B}\(300\) & \sk{C}\(450\) & \sk{D}\(600\) & \sk{E}\(900\)\end{tabular*}}
\end{minipage}
\par\vspace{\soruara}
\end{multicols}

\vspace{6pt}
\needspace{8\baselineskip}
\noindent\textbf{\color{anamavi}Cevap tablosu} {\small (işaretlediğiniz şıkları buraya da yazabilirsiniz)}\par\smallskip
\noindent\renewcommand{\arraystretch}{1.6}\begin{tabular}{|>{\centering\arraybackslash}p{1.25cm}|>{\centering\arraybackslash}p{1.25cm}|>{\centering\arraybackslash}p{1.25cm}|>{\centering\arraybackslash}p{1.25cm}|>{\centering\arraybackslash}p{1.25cm}|>{\centering\arraybackslash}p{1.25cm}|>{\centering\arraybackslash}p{1.25cm}|>{\centering\arraybackslash}p{1.25cm}|>{\centering\arraybackslash}p{1.25cm}|>{\centering\arraybackslash}p{1.25cm}|}\hline
\textbf{1} & \textbf{2} & \textbf{3} & \textbf{4} & \textbf{5} & \textbf{6} & \textbf{7} & \textbf{8} & \textbf{9} & \textbf{10}\\\hline
 &  &  &  &  &  &  &  &  & \\\hline
\textbf{11} & \textbf{12} & \textbf{13} & \textbf{14} & \textbf{15} & \textbf{16} & \textbf{17} & \textbf{18} & \textbf{19} & \textbf{20}\\\hline
 &  &  &  &  &  &  &  &  & \\\hline
\textbf{21} & \textbf{22} & \textbf{23} & \textbf{24} & \textbf{25} &  &  &  &  & \\\hline
 &  &  &  &  &  &  &  &  & \\\hline
\end{tabular}

\vfill\centering{\small Başarılar!}
\end{document}
